\documentclass[10pt]{amsart}
\usepackage[a4paper,margin=22mm]{geometry}
\usepackage[no-math]{fontspec}
\setmainfont{Latin Modern Roman}[SmallCapsFont={lmromancaps10-regular.otf}]
\usepackage{amsmath,amssymb,amsthm,mathtools,hyperref,graphicx,xfrac,etoolbox}
\hypersetup{hidelinks}
\usepackage[scr=rsfs,cal=euler]{mathalfa}
\newcommand{\ie}{i.e.}
\newcommand{\eg}{e.g.}
\newcommand{\cf}{cf.}
\newcommand{\resp}{resp.}
\newcommand{\skpt}{\par\smallskip\noindent}
\let\Subsection\subsection
\let\Subsubsection\subsubsection
\emergencystretch=3em
\newcommand\mto{\mathchoice{\longmapsto}{\mapsto}{\mapsto}{\mapsto}}
\newcommand{\Psfrac}[2]{(\sfrac{#1}{#2})}
\newcommand{\pspfrac}[2]{\sfrac{(#1)}{(#2)}}
\newcommand{\com}{\textup{com}}

\AtBeginDocument{%
\renewcommand{\labelitemi}{$\scriptscriptstyle\bullet$}
}

\let\ts\textstyle

\theoremstyle{plain}
\newtheorem{theorem}{Theorem}[section]
\newtheorem{lemma}[theorem]{Lemma}
\newtheorem{question}[theorem]{Question}
\newtheorem{corollary}[theorem]{Corollary}
\newtheorem{proposition}[theorem]{Proposition}

\theoremstyle{definition}
\newtheorem{definition}[theorem]{Definition}
\newtheorem{example}[theorem]{Example}
\newtheorem{remark}[theorem]{Remark}

\newcommand{\C}{\mathbf{C}}
\newcommand{\Z}{\mathbf{Z}}
\newcommand{\cX}{\mathcal{X}}
\newcommand{\cM}{\mathcal{M}}
\newcommand{\cN}{\mathcal{N}}
\newcommand{\E}{\mathbf{E}}
\newcommand{\vn}{\mathcal{C}}
\newcommand{\cH}{\mathcal{H}}
\newcommand{\cU}{\mathcal{U}}
\newcommand{\prob}{\mathbf{P}}
\newcommand{\F}{\mathbf{F}}
\newcommand{\anticom}{\textup{anticom}}

\DeclareMathOperator{\val}{val}
\DeclareMathOperator{\tr}{tr}
\DeclareMathOperator{\Tr}{Tr}
\numberwithin{equation}{section}
\title[HS stability of direct products]{Hilbert-Schmidt stability of direct products of countable stable groups when one factor has Kazhdan property (T)}
\author{Mikael De La Salle}
\date{Published version; DOI 10.5802/jep.314}
\begin{document}
\maketitle
\setcounter{section}{0}\setcounter{theorem}{0}
\begin{theorem}\label{thm:stability_direct_product_simple} Let $G$ and $H$ be two countable groups, one of which has property~(T).\enlargethispage{.5\baselineskip}

\begin{itemize}
\item
If $G$ and $H$ are both Hilbert-Schmidt stable, then so is $G\times H$.

\item
If $G$ and $H$ are both Hilbert-Schmidt flexibly stable, then so is $G\times H$.
\end{itemize}
\end{theorem}
\setcounter{section}{0}
\section{Preliminaries}\label{sec:preliminaries}
\subsection{Matrix and von Neumann algebras notation}
We recall a few basic facts about von Neumann algebras. For proofs and much more, we refer the reader to standard texts such as \cite{MR1873025}, or the modern book project \cite{AnantharamanPopa}.

A von Neumann algebra is a self-adjoint subalgebra of the algebra $B(\cH)$ of bounded operators on a complex Hilbert space $\cH$ that is equal to its bicommutant. Here the commutant of a subset $A\subset B(\cH)$ is the algebra
\[
A':=\{X\in B(\cH)\mid \forall Y \in A,\ XY=YX\},
\] and its bicommutant is the commutant of its commutant. Every von Neumann algebra is uniquely a dual space \cite[Cor.\,III.3.9]{MR1873025}, which allows us to talk about its weak-* topology.

A state on a von Neumann algebra $\cM$ is a linear form $\cM \to \C$ that is positive, meaning that $\tau(x^*x)\geq 0$ for every $x\in \cM$, and normalized by $\tau(1)=1$. It is called faithful if $\tau(x^*x)=0$ holds only if $x=0$. It is called normal if it is continuous for the weak-* topology, and it is called a trace if $\tau(ab)=\tau(ba)$.

In the whole note, we denote by $(\cM,\tau)$ a von Neumann algebra with a normal faithful tracial state. The main case of interest in $\cM = M_n(\C)$ for large $n$ and $\tau$ is the normalized trace $\tr = \frac{1}{n}\Tr$. The reader can safely assume throughout the paper that we are in this situation.

We will denote $\|x\|_2 = \big(\tau(x^*x)\big)^{\sfrac 1 2}$ for the $L_2$ norm on $\cM$. The completion of $\cM$ for this norm is denoted $L_2(\cM,\tau)$. We also use the notation $\|x\|_q = \big(\tau( (x^*x)^{q/2})\big)^{\sfrac 1 q}$ for the $L_q$ norm and $1\leq q<\infty$, and $\|x\|_\infty$ for the operator norm.

If $(\cM,\tau)$ is a von Neumann algebra with a faithful normal tracial state and if $\cN \subset\cM$ is an inclusion of von Neumann algebras, we obtain that $L_2(\cN,\tau)\subset L_2(\cM,\tau)$. The orthogonal projection $L_2(\cM,\tau) \to L_2(\cN,\tau)$ turns out to map $\cM$ into $\cN$, and the corresponding map is called the conditional expectation from $\cM$ to $\cN$ \cite[\S 9]{AnantharamanPopa} or \cite[Prop.\,V.2.36]{MR1873025}.

We also denote $\cM_\infty = \cM \overline{\otimes} B(\ell_2)$ with (infinite) trace $\tau_\infty:=\tau \otimes \Tr$. We often identify $\cM$ with $\cM \otimes e_{1,1} \subset \cM_\infty$ (with a different unit $1_\cM = 1 \otimes e_{1,1}$ than $1_{\cM_\infty}$), and often use the same letter $\tau$ to denote the amplified trace $\tau_\infty$. We also use the notation $\|x\|_2 \!=\! \big(\tau_\infty(x^*x)\big)^{\sfrac 1 2}$ for $x \!\in\! \cM_\infty$, but in this situation the norm can take the value~$\infty$.
\setcounter{subsection}{2}
\subsection{Spectral gap preliminaries}\label{sec:spectral_gap}
Let $G$ be a countable group (for example a finite group). If~$\mu$ is a symmetric probability measure on $G$ with generating support, then for every unitary representation $(\pi,\cH)$, $\pi(\mu):=\sum_g \mu(g)\pi(g)$ is a self-adjoint operator of norm $\leq 1$, whose eigenspace for the eigenvalue $1$ is the space of invariant vectors $\cH^\pi=\{\xi\in\cH \mid \forall g\in G, \pi(g)\xi=\xi\}$ (this uses that the support of $\mu$ generates the whole group $G$). We say that $\mu$ has spectral gap is there is a positive real number $\kappa$ such that for every unitary representation $(\pi,\cH)$ of $G$, the spectrum of $\pi(\mu)$ in contained in $[-1,1-\sfrac{1}{\kappa}]\cup\{1\}$, or equivalently is the restriction of $\pi(\mu)$ to the orthogonal of~$\cH^\pi$ has spectrum contained in $[-1,1-\sfrac{1}{\kappa}]$. The smallest such $\kappa$ is denoted by $\kappa(\mu)$, it~is the inverse of the spectral gap. If~$\mu$ does not have spectral gap, we set $\kappa(\mu)=\infty$. The reason for this parametrization of the spectral gap is because it is proportional to the constant appearing in the Poincaré inequalities, and is related to the Kazhdan constants. Specifically, $\kappa(\mu)$ is the smallest real number such that, for every unitary representation $(\pi,\cH)$ of $G$, and every vector $\xi \in \cH$,
\begin{equation}\label{eq:Poincare_constant} \|\xi - P_{\cH^\pi} \xi\|^2 \leq \frac\kappa 2 \int_G \|\pi(g) \xi-\xi\|_2^2\,d\mu(g).
\end{equation}
Here $P_{\cH^\pi}$ is the orthogonal projection on $\cH^\pi$. Let us justify this characterization of~$\kappa$. Indeed, by expanding the square norms, \eqref{eq:Poincare_constant} is equivalent to the inequality
\[
\|\xi\|^2 \leq \kappa (\|\xi\|^2 - \langle\pi(\mu)\xi,\xi\rangle)\quad \forall \xi \in (\cH^\pi)^\perp,
\]
or in other words (Rayleigh quotients) the spectrum of $\pi(\mu)$ restricted to $(\cH^\pi)^\perp = \ker(\pi(\mu)-1)^\perp$ consists of real numbers $\lambda$ satisfying $1\leq \kappa(1-\lambda)$, that is, $\lambda \leq 1-\sfrac{1}{\kappa}$.

For a given group $G$, $\mu$ has spectral gap if and only if $G$ has Kazhdan's property~(T). In~particular, the property that $\kappa(\mu)$ is finite does not depend on $\mu$. But it will be useful, in particular for finite groups, to study measures with good spectral gap. The extreme case is when $\mu =\prob_G$ is the uniform probability on $G$. In~that case, $\pi(\prob_G)$ is the orthogonal projection on $\cH^\pi$ and therefore $\kappa(\mu) =1$ (if $G$ is not the group $\{0\}$).

If $\mu$ is a not-necessarily symmetric probability measure on $G$ with support still generating, we define $\kappa(\mu)$ as $\kappa(\nu)$ where $\nu$ is the symmetric measure\vspace*{-3pt}
\[
\nu(g) = \frac{1}{2}(\mu(g)+\mu(g^{-1}).
\]
It is also the smallest constant such that \eqref{eq:Poincare_constant} holds.


\section*{Necessary spectral-gap commutator lemma (source Lemma 2.2)}
\setcounter{section}{2}\setcounter{theorem}{1}\setcounter{equation}{0}
\begin{lemma}\label{lemma:spectral_gap_commutator} Let $G$ be a group, $\mu$ a probability measure on $G$ with generating support, and $U\colon G\to \cU(\cM)$ be a group homomorphism. Set $\cN=\cM \cap U(G)'$ and let \hbox{$E_\cN:\cM\to\cN$} be the conditional expectation.

For every $V \in \cM$,
\begin{equation}\label{eq:spectral_gap_and_commutators} \| V - E_\cN(V)\|_2^2 \leq \frac{\kappa(\mu)}{2} \int \|[U(g),V]\|_2^2 \,d\mu(g).
\end{equation}

If $G$ is finite,
\[
\E_{g \in G} \|[U(g),V]\|_2^2 \leq \kappa(\mu) \int \|[U(g),V]\|_2^2 \,d\mu(g).
\]
\end{lemma}
\begin{proof} The formula
\[
\pi(a) X = U(a) X U(a)^*
\]
defines a unitary representation of $\pi$ on the Hilbert space $L_2(\cM,\tau)$. Here we use in an essential way that $\tau$ is trace. Moreover, the space of invariant vectors $L_2(\cM,\tau)^\pi$ coincides with $L_2(\cN,\tau)$, and the orthogonal projection on $L_2(\cM,\tau)^\pi$ is just $E_\cN$. Therefore, \eqref{eq:spectral_gap_and_commutators} is exactly the Poincaré inequality \eqref{eq:Poincare_constant}.

If $G$ is finite, $E_\cN$ is given by $V\mto \E_{g\in G} U(g) V U(g)^*$, so
\begin{align*} \E_{g \in G} \|[U(g),V]\|_2^2 &= 2\|V\|_2^2-2\tau(\E_{g \in G} U(g)^* V U(g)V^*)\\& = 2\|V\|_2^2-2\|E_\cN(V)\|_2^2 = 2\|V-E_\cN(V)\|_2^2.
\end{align*}
So the second inequality is a particular case of \eqref{eq:spectral_gap_and_commutators}.
\end{proof}
\setcounter{section}{3}
\section{Stability}\label{sec:stability}

\begin{definition} Let $G$ be a countable group and $\mu$ a probability measure with generating support. A map $\varphi:G\to \cU(\cM)$ is said to be an $(\varepsilon,\mu)$-almost homomorphism if $\iint \| \varphi(gh) - \varphi(g) \varphi(h)\|_2^2\,d\mu(g) \,d\mu(h) \leq \varepsilon$.
\end{definition}
Let $\vn$ be a class of von Neumann algebras equipped with normal tracial states.

\begin{definition}\label{def:stability}
We say that $(G,\mu)$ is $\vn$-stable if there is a non-decreasing
function $\delta:[0,4]\to [0,4]$ with $\lim_{t \to 0} \delta(t)= 0$ and such that, for every $(\cM,\tau)$ in $\vn$ and every $(\varepsilon,\mu)$-almost homomorphism $\varphi: G \to \cU(\cM)$, there is a group homomorphism $\pi:G \to
\cU(\cM)$ satisfying
\[
\int \|\varphi(g) - \pi(g)\|_2^2 \,d\mu(g) \leq \delta(\varepsilon).
\]

Such a function $\delta$ satisfying moreover that $\delta$ is concave and $\delta(4)= 4$ will be called a modulus of $\vn$-stability.
\end{definition}
\begin{remark} If $(G,\mu)$ is $\vn$-stable, then it admits a modulus of $\vn$-stability (by~replacing $\delta$ by the smallest function greater than $\delta$, concave and taking the value~$4$ at~$4$). The concavity requirement for $\delta$ is here to make statements such as Theorem~\ref{thm:stability_direct_product} cleaner. It is also very natural, and in fact if $\vn$ is stable by direct sums then the best $\delta$ is necessarily concave. Indeed, if $\varphi_1$ and $\varphi_2$ are respectively $(\varepsilon_1,\mu)$ and $(\varepsilon_2,\mu)$ almost representations with values in $\cU(\cM_1)$ and $\cU(\cM_2)$, and $\lambda \in [0,1]$, then defining $\cM = \cM_1\oplus \cM_2$ with trace $\tau(x_1,x_2) = \lambda \tau_1(x_1) + (1-\lambda)\tau_2(x_2)$, the pair $(\varphi_1,\varphi_2)$ defines a $(\lambda\varepsilon_1+(1-\lambda)\varepsilon_2,\mu)$ almost representation $\varphi$ on $\cM$. Moreover, a unitary representation $\pi:G\to \cU(\cM)$ is the same as a pair of unitary representations $\pi_i:G\to \cU(\cM_i)$, with
\begin{multline*} \int \|\varphi(g) - \pi(g)\|_2^2 \,d\mu(g) \\=\lambda \int \|\varphi_1(g) - \pi_1(g)\|_2^2 \,d\mu(g) + (1-\lambda) \int \|\varphi_2(g) - \pi_2(g)\|_2^2 \,d\mu(g).
\end{multline*}
Taking the supremum over all $(\varepsilon_1,\mu)$ and $(\varepsilon_2,\mu)$ almost representations, we obtain
\[
\lambda\delta(\varepsilon_1) + (1-\lambda)\delta(\varepsilon_2) \leq \delta(\lambda\varepsilon_1+(1-\lambda)\varepsilon_2).
\]
That is, $\delta$ is concave.
\end{remark}

\begin{definition}\label{def:flexible_stability}
We say that $(G,\mu)$ is $\vn$ flexibly stable if there is a non-decreasing function $\delta:[0,4]\to [0,4]$ such that $\lim_{t \to 0} \delta(t) = 0$ such that, for every $(\cM,\tau)$ in $\vn$ and every $(\varepsilon,\mu)$-almost homomorphism $\varphi: G \to \cU(\cM)$, there is a projection $P \in \cM_\infty$, a group homomorphism $\pi:G \to \cU(\cM)$ and an isometry $w \in P \cM_\infty 1_\cM$ satisfying
\[
\max\left(\tau(P)-1, \int \|\varphi(g) - w^*\pi(g)w\|_2^2 \,d\mu(g)\right) \leq \delta(\varepsilon).
\]

Such a function $\delta$ satisfying moreover that $\delta$ is concave and $\delta(4)=4$ will be called a modulus of $\vn$-flexible stability.
\end{definition}
This is an adaptation of the standard notions of Hilbert-Schmidt stability and Hilbert-Schmidt flexible stability, see for example \cite{Ioana2}.

Let $\vn_{\mathrm{fin}}$ denote the class of all finite-dimensional von Neumann algebras with tracial states.
\begin{lemma} A countable group $G$ is Hilbert-Schmidt (flexibly) stable if and only if there is a probability measure $\mu$ on $G$ with generating support such that $(G,\mu)$ is $\vn_{\mathrm{fin}}$ (flexibly) stable.

If $G$ is finitely presented, $\mu$ can moreover be taken to be of finite support.
\end{lemma}
\begin{proof}
The lemma is immediate if $\vn_{\mathrm{fin}}$ is replaced by the set \[
\{M_n\} = \{(M_n(\C),\tr) \mid n\geq 1\}.
\] So the whole point of the lemma is to show that if $(G,\mu)$ is $\{M_n\}$-stable, then $(G,\mu)$ is $\vn_{\mathrm{fin}}$-stable, and similarly for flexible stability.

Assume that $(G,\mu)$ is $\{M_n\}$-stable, and let $\delta$ be a modulus. We shall prove that $(G,\mu)$ is $\vn_{\mathrm{fin}}$-stable with the same modulus. Let $(\cM,\tau) \in \vn_{\mathrm{fin}}$ and $\varphi:G\to \cU(\cM)$ be an $(\varepsilon,\mu)$-almost homomorphism. We can decompose $\cM$ into a finite direct sum of matrix algebras $\cM = \bigoplus_{i=1}^k M_{n_i}(\C)$ with trace $\tau = \sum_i \lambda_i \tr_{n_i}$ for positive numbers $\lambda_i$ summing to $1$. Then $\varphi(g) = (\varphi_i(g))_i$ where $\varphi_i$ is an $(\varepsilon_i,\mu)$-almost homomorphism for some numbers $\varepsilon_i$ satisfying
\[
\sum\nolimits_i \lambda_i\varepsilon_i\leq\varepsilon.
\]
Therefore, by $\{M_n\}$-stability, there are homomorphisms $\pi_i \colon G\to\cU(n_i)$ such that $\int\|\varphi_i(g) - \pi_i(g)\|_2^2\,d\mu(g) \leq \delta(\varepsilon_i)$. If~we define a homomorphism $\pi:G\to \cU(\cM)$ by $\pi(g) = (\pi_i(g))_{g\in G}$, we obtain
\[
\int \|\varphi(g) - \pi(g)\|_2^2 \,d\mu(g) = \sum_{i=1}^k\lambda_i \int\|\varphi_i(g) - \pi_i(g)\|_2^2\,d\mu(g) \leq \delta(\varepsilon)
\]
by concavity of $\delta$.

The argument is identical for flexible stability, and left to the reader.
\end{proof}

Assume that $\vn$ is a class of von Neumann algebras closed under taking subalgebras.

\begin{theorem}\label{thm:stability_direct_product} Direct products of $\vn$-stable groups are $\vn$-stable, provided that one of them has property (T).

More precisely, if $(G_1,\mu_1)$ has property (T) and is $\vn$-stable with modulus $\delta_1$ and $(G_2,\mu_2)$ is $\vn$-stable with modulus $\delta_2$, then $(G_1 \times G_2,\mu)$ is $\vn$-stable with modulus $\delta(\varepsilon) \lesssim \delta_2(\kappa(\mu_1) \delta_2(\varepsilon))$, where $\mu$ is the probability measure
\[
\mu(x,y) = \frac{1}{2}(\mu(x) 1_{y=1}+\mu(y) 1_{x=1}).
\]
\end{theorem}
The proof will use the following simple lemmas.
\begin{lemma}\label{eq:almost-unitary_close_to_unitary} Let $(\cM,\tau)$ be a tracial von Neumann algebra and $\cN\subset\cM$ be a subalgebra, and let $E_\cN:\cM \to \cN$ be the conditional expectation. For every unitary $V \in \cU(\cM)$, there is a unitary $\tilde{V}\in\cU(\cN)$ such that $\|V-\tilde{V}\|_2 \leq \sqrt{2}\, \|V-E_\cN(V)\|_2$.
\end{lemma}
The $\sqrt{2}$ is optimal, for example when $E_\cN(V)=0$.
\begin{proof} Let $X = E_\cN(V)$. Since $\cN$ is finite, we can write $X=\tilde{V}|X|$ where $\tilde{V} \in\cU(\cN)$ and $|X|=(X^*X)^{\sfrac 1 2}$. By the duality between $L_1(\cM,\tau)$ and $\cM$, we have $\Re \tau(V^*X) \leq\tau(|X|)$, or equivalently $\|\tilde{V}-X\|_2 \leq \|V-X\|_2$. If~we decompose $V-\tilde{V} = V-X + X-\tilde{V}$, the terms $V-X$ and $X-\tilde{V}$ are orthogonal, and therefore
\[
\|V-\tilde{V}\|_2^2 = \|V-X\|^2 + \|\tilde{V}-X\|^2 \leq 2\,\|V-X\|_2^2.
\]
This proves the lemma.
\end{proof}
\begin{lemma}\label{eq:norm_of_conditional_exp} Let $\cN\subset\cM$ be a von Neumann subalgebra, and $E_\cN:\cM\to\cN$ be the trace-preserving conditional expectation. For every $\xi \in L_2(\cM,\tau)$,
\[
\|\xi - E_\cN(\xi)\|_2 = \sup\{\tau(\xi \eta) \mid \eta \in L_2(\cM,\tau),\, E_\cN(\eta)=0,\,\|\eta\|_2 = 1\}.
\]
\end{lemma}
\begin{proof}
$(1-E_\cN)$ is the orthogonal projection on $\{\eta \in L_2(\cM,\tau) \mid E_\cN(\eta)=0\}$.
\end{proof}

\begin{proof}[Proof of Theorem~\ref{thm:stability_direct_product}]
Let $\varphi: G_1\times G_2\to\cU(\cM)$ be a $(\varepsilon,\mu)$-almost homomorphism. The idea is simple: we first use the stability of $G_1$ to say that the restriction of $\pi$ to $G_1$ is close to a representation. Then by property (T), we will deduce that the restriction of $\pi$ to $G_2$ is close to an almost homomorphism from $G_2$ to the commutant of $G_1$, and therefore by stability of $G_2$ to an actual homomorphism with values in the commutant of $G_1$.

Here are the details. If~we denote
\[
\varepsilon_{i,j} = \iint_{G_i\times G_j} \| \varphi(gh) - \varphi(g) \varphi(h)\|_2^2\,d\mu_i(g) \,d\mu_j(h),
\]
we then have $\varepsilon_{1,1}+\varepsilon_{2,2} + \varepsilon_{1,2}+\varepsilon_{2,1}\leq 4\varepsilon$.

The restriction of $\varphi$ to $G_1$ is an $(\varepsilon_{1,1},\mu_1)$-almost representation of $G_1$, so there is a unitary representation $\pi_1:G_1\to\cU(\cM)$ such that
\[
\int \|\varphi(g)-\pi_1(g)\|_2^2 \,d\mu_1(g) \leq \delta_1(\varepsilon_{1,1}).
\]
Set $\cN=\cM \cap \tilde{\pi}_1(G_1)'$, and let $E:\cM\to\cN$ be the conditional expectation. Since $\vn$ is assumed to be stable under taking subalgebras, we have $\cN\in\vn$. Define, for every $h \in G_2$,
\[
\eta(h) = \|\varphi(h) -E(\varphi(h))\|_2.
\]
We first establish an upper bound for the norm of $\eta$ in $L_2(G_2,\mu_2)$. This is where we use that $G_1$ has property (T). By Lemma~\ref{lemma:spectral_gap_commutator}, we know that
\[
\eta(h)^2 \leq \kappa(\mu_1) \int_{G_1} \|[\varphi(h),\pi_1(g)]\|_2^2 \,d\mu_1(g).
\]
We can bound
\[
\|[\varphi(h),\pi_1(g)\|_2 \leq 2\|\varphi(g)-\pi_1(g)\|_2 + \|\varphi(h)\varphi(g) -\varphi(hg)\|_2 + \|\varphi(g) \varphi(h) - \varphi(gh)\|_2.
\]
So, by the triangle inequality in $L_2(\mu_1\times \mu_2)$, we obtain
\[
\|\eta\|_{L_2(\mu_2)} \leq \sqrt{\kappa(\mu_1)/2}\,\Bigl(2 \sqrt{\delta_1(\varepsilon_{1,1})} + \sqrt{\varepsilon_{1,2}} + \sqrt{\varepsilon_{2,1}}\Bigr).
\]
By the standing assumptions on $\delta_1$, and the Cauchy-Schwarz inequality, we obtain the following inequality that we will use shortly
\begin{equation}\label{eq:L2norm_of_eta} \int\eta(h)^2 \,d\mu_2(h) \leq 12 \kappa(\mu_1) \delta_1(\varepsilon).\end{equation}

By Lemma~\ref{eq:almost-unitary_close_to_unitary}, for every $h \in G_2$ there is a unitary $V(h) \in \cU(\cN)$ such that $\|\varphi(h) -\nobreak V(h)\|_2\leq \sqrt{2}\eta(h)$. Our goal is to show that $V$ is an almost-homomorphism, to then apply flexible stability for $G_2$. By the triangle inequality,
\[
\left(\int \|V(gh)-V(g)V(h)\|_2^2 \,d\mu_2(g) \,d\mu_2(h)\right)^{\sfrac 1 2}
\]
is bounded above by
\[
\sqrt{\varepsilon_{2,2}} + 2\sqrt{2}\, \|\eta\|_{L_2(\mu_2)} + \sqrt{2}\,\|\eta\|_{L_2(\mu_2\ast\mu_2)}.
\]
Decompose $\varphi(gh) = a+b+c$ with $a= \varphi(gh)-\varphi(g)\varphi(h)$, $b= \varphi(g) (\varphi(h) - E_\cN(\varphi(h)))$ and $c=\varphi(g) E_\cN(\varphi(h))$. Using Lemma~\ref{eq:norm_of_conditional_exp}, and writing $\sup_\eta$ for the supremum over all unit vectors in the orthogonal subspace of $L_2(\cN,\tau)$ in $L_2(\cM,\tau)$, we have
\begin{align*}
\eta(gh) & = \sup\nolimits_\eta \tau(a\eta) + \tau(b\eta)+ \tau(c\eta)\\
& \leq \|a\|_2 + \|b\|_2 +\|c - E_\cN(c)\|_2\\
& \leq \|\varphi(gh) - \varphi(g) \varphi(h)\|_2 + \eta(h)+\eta(g).
\end{align*}
As a consequence, we have $\|\eta\|_{L_2(\mu_2\ast\mu_2)} \leq \sqrt{\varepsilon_{2,2}} + 2\,\|\eta\|_{L_2(\mu_2)}$, and we deduce
\[
\left(\int \|V(gh)-V(g)V(h)\|_2^2 \,d\mu_2(g) \,d\mu_2(h)\right)^{\sfrac 1 2} \leq (1+\sqrt{2})\sqrt{\varepsilon_{2,2}} + 4\sqrt{2}\, \|\eta\|_{L_2(\mu_2)}.
\]
By \eqref{eq:L2norm_of_eta}, this last quantity is less than $\sqrt{C\kappa(\mu_1) \delta_1(\varepsilon)}$ for some universal constant~$C$. The theorem follows easily. Indeed, by stability for $G_2$, we deduce that there is a unitary representation $\pi_2:G\to\cU(\cN)$ such that
\[
\int \|V(g) - \pi_2(g)\|_2^2\,d\mu_2(g) \leq \delta_2(C\kappa(\mu_1)\delta_1(\varepsilon)),
\]
and therefore
\[
\int \|\varphi(g)- \pi_2(g)\|_2^2\,d\mu_2(g) \leq 3\bigl(\|\eta\|_{L_2(\mu_2)}^2 + \delta_2(C\kappa(\mu_1)\delta_1(\varepsilon))\bigr).
\]
By definition of $\cN$, $\pi_2(g)$ commutes with $\pi_1(G_1)$, so the pair $({\pi}_1,{\pi}_2)$ gives rise to a unitary representation ${\pi}: G_1\times G_2\to \cU(\cM)$ by ${\pi}(g,h)={\pi}_1(g){\pi}_2(h)$. We have
\[
\int \|\varphi(g) - {\pi}(g)\|_2^2 \,d\mu(g)
\leq \frac{1}{2}\bigl(\delta_1(\varepsilon_{1,1}) + 3\|\eta\|_{L_2(\mu_2)}^2 + 3\delta_2(C\kappa(\mu_1)\delta_1(\varepsilon))\bigr).
\]
This is $\lesssim \delta_2(\kappa(\mu_1)\delta_1(\varepsilon))$ by \eqref{eq:L2norm_of_eta} and the standing assumption that the moduli $\delta_i$ are concave and satisfy $\delta_i(t) \geq t$.
\end{proof}
\begin{theorem}\label{thm:flexible_stability_direct_product} Direct products of $\vn$-flexibly stable groups are $\vn$-flexibly stable, provided that one of them has property (T).

More precisely, if $(G_1,\mu_1)$ has property (T) and is $\vn$-flexibly stable with modulus $\delta_1$ and $(G_2,\mu_2)$ is $\vn$-flexibly stable with modulus $\delta_2$, then $(G_1 \times G_2,\frac{1}{2}(\mu_1 \otimes \delta_1 + \delta_1 \otimes \mu_2))$ is $\vn$-flexibly stable with modulus $\delta(\varepsilon) \lesssim \delta_2(\kappa(\mu_1) \delta_2(\varepsilon))$.
\end{theorem}
\begin{proof}
This is proved in the same way as for stability. The only difference is that the von Neumann algebras change after each use of flexible stability. First, when flexible stability is used for $G_1$, a representation $\tilde{\pi}_1$ of $G_1$ is constructed in a small dilation $\cM_1$ of $\cM$. The almost representation of $G_2$ is then defined with values in $\cU(\cM_1)$ as $w_1 \pi(g) w_1^* +P_1 - w_1w_1^*$. Similarly, when flexible stability is used for $G_2$ to construct a representation $\tilde{\pi}_2$ in a small dilation $\cM_2$ of $\cM_1$, the representation $\tilde{\pi}_1$ of~$G_1$ is then defined with values in $\cU(\cM_2)$ as
$w_2 \tilde{\pi}_1(g) w_2^* +P_2 - w_2w_2^*$.
\end{proof}
\begin{thebibliography}{99}
\bibitem[2]{AnantharamanPopa} C. Anantharaman \& S. Popa - An introduction to II 1 factors, Book available https://www.idpoisson.fr/anantharaman/publications/IIun.pdf.
\bibitem[10]{Ioana2} A. Ioana - “Almost commuting matrices and stability for product groups”, J. Eur. Math. Soc. (JEMS) 27 (2025) no. 10, p. 4027-4068.
\bibitem[20]{MR1873025} M. Takesaki - Theory of operator algebras. I, Encyclopaedia of Math. Sciences, vol. 124, Springer-Verlag, Berlin, 2002.
\end{thebibliography}
\end{document}
